\documentclass[5p,times,twocolumn]{elsarticle}

\journal{SoftwareX}

\usepackage{amsmath}
\usepackage{amssymb}
\usepackage{booktabs}
\usepackage{tabularx}
\usepackage{listings}
\usepackage{xcolor}
\usepackage[hidelinks]{hyperref}

\lstset{
  basicstyle=\ttfamily\scriptsize,
  keywordstyle=\color{blue!70},
  commentstyle=\color{gray},
  stringstyle=\color{orange!80!black},
  language=Python,
  breaklines=true,
  breakatwhitespace=true,
  frame=single,
  numbers=left,
  numberstyle=\tiny\color{gray},
  xleftmargin=1.5em,
  framexleftmargin=1.5em,
}

\begin{document}

\begin{frontmatter}

\title{MOLEKUL: A Pure-Python Ab Initio Quantum Chemistry Platform for
       Education and Reproducible Benchmarking}

\author{Enis Yazici\corref{cor1}}
\ead{yazicienis@gmail.com}
\cortext[cor1]{Corresponding author}
\address{SRH University of Applied Sciences Heidelberg,
         School of Engineering and Architecture,
         Heidelberg, Germany}

\begin{abstract}
MOLEKUL is an open-source, pure-Python implementation of Restricted
Hartree--Fock (RHF) self-consistent field theory and second-order
M{\o}ller--Plesset perturbation theory (MP2) for closed-shell molecules.
Starting from a molecular geometry and a contracted Gaussian basis set, MOLEKUL
evaluates all one- and two-electron integrals analytically, solves the
Roothaan--Hall equations iteratively, and reports the total RHF and MP2
energies.
Built-in basis data cover hydrogen through fluorine for STO-3G, 6-31G*, and
cc-pVDZ.
Additional modules provide a geometry optimizer and Mulliken and L{\"o}wdin
population analysis.
Validated against PySCF on a 14-molecule benchmark suite, MOLEKUL achieves
RHF/STO-3G total-energy agreement within $5\times10^{-8}$~E$_h$ and MP2
correlation-energy agreement within $2\times10^{-7}$~E$_h$.
The codebase comprises 606 automated tests and is designed for advanced
undergraduate and graduate courses in quantum chemistry, as well as for
reproducible NumPy/BLAS performance profiling.
\end{abstract}

\begin{keyword}
quantum chemistry \sep
Hartree--Fock \sep
MP2 \sep
ab initio \sep
Gaussian basis sets \sep
Python \sep
education
\end{keyword}

\end{frontmatter}

% ============================================================
\section*{Code Metadata}
% ============================================================

\begin{table*}[tbp]
\caption{Code metadata}
\label{tab:metadata}
\begin{tabularx}{\textwidth}{lX}
\toprule
\textbf{Nr.} & \textbf{Detail} \\
\midrule
C1 & Current code version: v0.1.1 \\
C2 & Repository: \url{https://github.com/yazicienis/molekul} \\
C3 & Reproducible capsule (Zenodo): \url{https://doi.org/10.5281/zenodo.19743306} \\
C4 & License: MIT \\
C5 & Versioning system: git \\
C6 & Languages and tools: Python~$\ge$~3.10, NumPy \\
C7 & Dependencies: \texttt{numpy}; \texttt{pytest} (development) \\
C8 & Documentation: \url{https://github.com/yazicienis/molekul} \\
C9 & Support: \texttt{yazicienis@gmail.com} \\
\bottomrule
\end{tabularx}
\end{table*}

% ============================================================
\section{Motivation and Significance}
% ============================================================

Production electronic-structure packages such as PySCF~\cite{sun2020},
Psi4~\cite{smith2020}, and ORCA~\cite{neese2020} provide broad functionality
and optimized implementations, but their internal complexity can obscure the
connection between textbook equations and executable algorithms.
Earlier educational codes such as PyQuante~\cite{quante2004} demonstrated the
value of readable implementations, but they are no longer as actively maintained
or systematically validated against modern reference packages.

MOLEKUL occupies the space between textbook equations and production code.
Every algorithmic step traces directly to a named Python function with clear
correspondence to standard quantum-chemistry
references~\cite{szabo1989,helgaker2000}.
The 606-test suite and deterministic NumPy implementation ensure that every
numerical claim is automatically reproducible.

Beyond pedagogy, the Python/NumPy implementation provides a self-contained
baseline for performance profiling, because quantum-chemistry algorithms are
implemented explicitly rather than delegated to compiled electronic-structure
kernels, allowing clean isolation of NumPy/BLAS performance characteristics.

% ============================================================
\section{Software Description}
% ============================================================

\subsection{Architecture}

MOLEKUL is organized as an installable Python package (\texttt{src/molekul/}).
Table~\ref{tab:modules} lists the main modules.
Installation: \texttt{pip install -e ".[dev]"}.
Tests: \texttt{pytest tests/}.

\begin{table}[tbp]
\caption{Main modules of MOLEKUL.}
\label{tab:modules}
\begin{tabularx}{\columnwidth}{lX}
\toprule
\textbf{Module} & \textbf{Responsibility} \\
\midrule
\texttt{integrals.py}   & One-electron integrals \\
\texttt{eri.py}         & Two-electron repulsion integrals \\
\texttt{rhf.py}         & RHF SCF driver \\
\texttt{mp2.py}         & MP2 correlation energy \\
\texttt{population.py}  & Mulliken and L{\"o}wdin analysis \\
\texttt{optimizer.py}   & Geometry optimizer \\
\texttt{harmonic.py}    & Harmonic frequencies (experimental) \\
\bottomrule
\end{tabularx}
\end{table}

\subsection{Integral Engine}

One- and two-electron integrals are evaluated analytically using the
McMurchie--Davidson scheme~\cite{mcmurchie1978,helgaker2000}, reducing all
integral evaluation to a Boys-function call
$F_n(x)=\int_0^1 t^{2n}e^{-xt^2}\mathrm{d}t$~\cite{boys1950}.

The Boys function is evaluated using a pure-Python combination of
\texttt{math.erf} for the analytic $F_0$ expression, recurrence relations, and
small-$x$ Taylor expansions, with the branch selected to maintain numerical
stability.
A dedicated test confirms relative errors below $5\times10^{-14}$ across the
argument range encountered in STO-3G through cc-pVDZ bases; no SciPy
dependency is required.

ERIs are stored in a full $N_\mathrm{AO}^4$ double-precision array---dense
storage is intentional, keeping the Fock-build loop transparent.
At $N_\mathrm{AO}=100$ this requires approximately 800~MB; practical use is
limited to small molecules.
Contracted shells are explicitly renormalized to unit overlap, correcting for
cross-primitive contributions commonly omitted in introductory treatments.

\subsection{SCF Driver}

The RHF SCF driver~\cite{roothaan1951,hall1951} implements:
(i)~a Superposition of Atomic Densities (SAD) initial guess with
spherically-averaged neutral-atom occupations;
(ii)~DIIS acceleration~\cite{pulay1980,pulay1982} using the error vector
$\mathbf{e}=\mathbf{FPS}-\mathbf{SPF}$; and
(iii)~an optional level shift~\cite{saunders1973} that stabilizes early
iterations.
Default thresholds: $|\Delta E|<10^{-10}$~E$_h$,
$\max|\Delta\mathbf{P}|<10^{-8}$.

\subsection{MP2 and Analysis}

The MP2~\cite{moller1934} correlation energy is computed via a four-index
integral transformation.
Mulliken and L{\"o}wdin~\cite{lowdin1950} population analysis and electric
dipole moments are also implemented.

\subsection{Basis Sets}

Built-in data (Table~\ref{tab:basis}) cover H--F; exponents and coefficients
are taken directly from PySCF's basis-set data~\cite{pritchard2019}.

\begin{table}[tbp]
\caption{Built-in basis sets.}
\label{tab:basis}
\begin{tabularx}{\columnwidth}{llX}
\toprule
\textbf{Basis} & \textbf{Reference} & \textbf{AO / heavy atom} \\
\midrule
STO-3G  & \cite{hehre1969}                          & 5  \\
6-31G*  & \cite{hehre1972,hariharan1973}            & 15 \\
cc-pVDZ & \cite{dunning1989}                        & 15 \\
\bottomrule
\end{tabularx}
\end{table}

% ============================================================
\section{Illustrative Example}
% ============================================================

The following script computes RHF and MP2 energies of water:

\begin{lstlisting}[language=Python]
from molekul.molecule import Molecule
from molekul.atoms import Atom
from molekul.basis_sto3g import get_sto3g
from molekul.rhf import rhf_scf
from molekul.mp2 import mp2_energy
import numpy as np

ANG2BOHR = 1.8897259886
mol = Molecule([
    Atom('O', np.array([ 0.  ,  0.   , 0.117])*ANG2BOHR),
    Atom('H', np.array([ 0.  ,  0.757,-0.469])*ANG2BOHR),
    Atom('H', np.array([ 0.  , -0.757,-0.469])*ANG2BOHR),
])
basis = get_sto3g()
rhf = rhf_scf(mol, basis)
mp2 = mp2_energy(mol, basis, rhf)
print(f"RHF: {rhf.energy_total:.8f} Eh")  # -74.96258854
print(f"MP2: {mp2.energy_total:.8f} Eh")  # -74.99844967
\end{lstlisting}

% ============================================================
\section{Validation}
% ============================================================

\subsection{RHF/STO-3G}

The script \texttt{scripts/benchmark\_14mol.py} compares MOLEKUL RHF/STO-3G
energies against PySCF~\cite{sun2020} for 14 closed-shell molecules (H$_2$
through CH$_2$O).
All 14 pass with maximum deviation $4.9\times10^{-8}$~E$_h$ (C$_2$H$_2$),
consistent with floating-point differences in integral evaluation.
Full results are in \texttt{outputs/logs/benchmark\_14mol.json}.

\subsection{MP2 Correlation Energy}

Table~\ref{tab:mp2} compares MP2 correlation energies against PySCF for eight
STO-3G molecules; all agree within $2\times10^{-7}$~E$_h$.

\begin{table}[tbp]
\caption{MP2 correlation energies: MOLEKUL vs.\ PySCF (STO-3G).}
\label{tab:mp2}
\begin{tabularx}{\columnwidth}{lXr}
\toprule
\textbf{Molecule} &
\textbf{$E_\mathrm{corr}^\mathrm{PySCF}$ / E$_h$} &
\textbf{$|\Delta E|$ / E$_h$} \\
\midrule
H$_2$  & $-0.01313807$ & $8\times10^{-10}$ \\
HF     & $-0.01734432$ & $3\times10^{-9}$  \\
H$_2$O & $-0.03550283$ & $1\times10^{-8}$  \\
N$_2$  & $-0.15419856$ & $2\times10^{-8}$  \\
CO     & $-0.12852242$ & $1\times10^{-7}$  \\
NH$_3$ & $-0.04704176$ & $6\times10^{-9}$  \\
CH$_4$ & $-0.05650741$ & $6\times10^{-9}$  \\
HCN    & $-0.12997898$ & $5\times10^{-9}$  \\
\bottomrule
\end{tabularx}
\end{table}

% ============================================================
\section{Known Limitations}
% ============================================================

Dense $N_\mathrm{AO}^4$ ERI storage limits practical use to small molecules
($N_\mathrm{AO}\lesssim50$); no integral screening, density fitting, or
integral-direct algorithms are implemented.
The SCF driver supports closed-shell RHF only (no UHF or ROHF).
Built-in basis data cover H--F; heavier elements are not parametrized.
Geometry optimization and harmonic-frequency workflows use finite differences
rather than analytic gradients.
Excited-state (CIS) and harmonic-frequency modules are experimental
pedagogical infrastructure and are not production-validated.
No effective core potentials or relativistic corrections are included.

% ============================================================
\section{Impact}
% ============================================================

MOLEKUL serves two complementary purposes.
As a teaching tool, all algorithms trace to named functions in readable Python,
making it suitable for advanced undergraduate and graduate courses in quantum
chemistry and computational physics.
As a performance-profiling baseline, the 606-test suite and deterministic
NumPy~\cite{harris2020} implementation allow comparisons of hardware and
numerical back-ends (e.g.\ Numba, CuPy) without compiled electronic-structure
kernels.
Benchmark timings and energies are logged to version-controlled JSON files.

% ============================================================
\section*{Software Availability}
% ============================================================

\begin{itemize}
  \item Repository: \url{https://github.com/yazicienis/molekul}
  \item Archive DOI: \url{https://doi.org/10.5281/zenodo.19743306}
  \item License: MIT
  \item Version archived: v0.1.1
\end{itemize}

% ============================================================
\section*{Declaration of Competing Interest}
% ============================================================

The author declares that he has no known competing financial interests or
personal relationships that could have appeared to influence the work reported
in this paper.

% ============================================================
\section*{Funding}
% ============================================================

This research did not receive any specific grant from funding agencies in the
public, commercial, or not-for-profit sectors.

% ============================================================
\section*{CRediT Author Statement}
% ============================================================

\textbf{Enis Yazici:} Conceptualization, Methodology, Software, Validation,
Investigation, Writing~-- original draft, Writing~-- review \& editing.

% ============================================================
\section*{Data Availability}
% ============================================================

The source code, validation scripts, and benchmark logs are available in the
public GitHub repository and archived on Zenodo at the DOI given in the Code
Metadata table.

% ============================================================
\section*{AI Usage Disclosure}
% ============================================================

Generative AI tools, including Claude (Anthropic) and ChatGPT (OpenAI), were
used for code drafting, refactoring, documentation support, test-generation
suggestions, and editorial review.
All scientific algorithms, numerical reference values, validation criteria, and
reported results were reviewed and verified by the author, who remains
responsible for the correctness, originality, and scientific content of the
submission.

% ============================================================
\section*{Acknowledgements}
% ============================================================

The author thanks the developers of PySCF, whose reference calculations were
used throughout validation, and the EMSL Basis Set Exchange~\cite{pritchard2019}
for tabulated basis-set parameters.

% ============================================================
\bibliographystyle{elsarticle-num}
\bibliography{paper}
% ============================================================

\end{document}
